% !TeX root = ../../protocolo.tex

 During my doctoral studies, I contributed to the publication of the following peer-reviewed work:
 
	\begin{tcolorbox}[boxsep = .75em, leftupper = 2.5em,			
		title = Phase Singularity in Random Gold Metasurface for Refractive Index Sensing: Experimental Demonstration and Theoretical Analysis \cite{cuanalo_phase_2024}]
%		
		J. P. Cuanalo-Fernández, N. Korneev, I. Cosme, M. B. De La Mora, \emphred{J. A. Urrutia-Anguiano}, A. Reyes-Coronado, R. Ramos-García, and S. Mansurova\\[.5em]
		\textit{J. Phys. Chem. C} \emphblue{128}(27): 11372–11381, 2024.\\[.5em]
		DOI: \href{https://pubs.acs.org/doi/10.1021/acs.jpcc.4c02274}{10.1021/acs.jpcc.4c02274}
	\end{tcolorbox}
%
	The phenomenon of topological darkness has been experimentally observed and  theoretically characterized in metasurfaces with an ordered spatial distributions of their building entities and exploited as a high sensitive sensing mechanism. In this work, the conditions to obtain topological darkness ---near zero reflectance and phase singularity--- are fulfilled in a random array of gold nanoislands (oblate ellipsoids) using theoretical and experimental methods. Reflectance and phase spectra are measured using a common-path interferometer in an attenuated total reflectance setup. Results reveal a partial state of topological darkness, with phase singularities marked by abrupt $\pm\pi$ phase jumps. Additionally, the study highlights the potential of phase and derivative measurements near the singularity for highly sensitive refractive index detection. The findings align well with theoretical predictions based on the Thin Film Theory.

Within the collaborations, I am working on the calculations and writing process of two manuscripts for submission as peer-reviewed publications:
%
\begin{itemize}
	\renewcommand\labelitemi{\textcolor{pblueFC}{$\blacktriangleright$}}	
	\item \emphblue{Title:} Theoretical and numerical approach to substrate effects in random plasmonic metasurfaces\\
			\emphblue{Authors:} \emphred{Jonathan A. Urrutia-Anguiano}, Isabel Y. Rojas-Martinez, Juan P. Cuanalo-Fernández, Alejandro Reyes-Coronado, Rubén Ramos-García, and %
							Svetlana Mansurova\\
			\emphblue{Status:} Review between collaborators\\
			\emphblue{Description:} When analyzing the optical response of metasurfaces, discrepancies between experimental measurements and theoretical predictions are often qualitatively attributed to the influence of the substrate. In this study, we compare experimentally obtained reflectivity spectra of disordered metasurfaces composed of spherical gold nanoparticles (AuNPs) with predictions from the Coherent Scattering Model (CSM). Deviations from the experimental results are further examined through Finite Element Method (FEM) simulations of a single AuNP. By quantitatively estimating the contribution of multiple scattering and of image multipoles induced in the substrate, our findings, supported by experimental data, FEM simulations, and CSM predictions suggest that the gold nanoparticles are partially embedded in the substrate. This partial embedding effectively alters the apparent size of the nanoparticles and increases the surface coverage of the metasurface.
%	
	\item \emphblue{Title:} Spectral shift in the strong couple regime: Red- and blueshift in plasmonic disordered metasurfaces\\
			\emphblue{Authors:}  Juan P. Cuanalo-Fernández, \emphred{Jonathan A. Urrutia-Anguiano}, Irving Gazga-Gurrión,  Nikolai Korneev, Alejandro Reyes-Coronado, Rubén Ramos-García, and %
			Svetlana Mansurova\\
			\emphblue{Status:} Writting process\\
			\emphblue{Description:} The optical response of metasurfaces composed of disordered gold nanoparticles can be modeled using multiple scattering theories, such as the Coherent Scattering Model and the Thin Film Theory. Under Attenuated Total Reflection conditions, the resonance of these metasurfaces exhibited a spectral shift in response to changes in the refractive index of the surrounding medium of the nanoparticles. In addition to the well-known redshift, both theoretical predictions and experimental observations have also revealed a blueshift. This blueshift is believed  to be an effect of strong coupling, as it occurs only when the angle of incidence is near the critical angle. Additionally, from a theoretical framework, it is studied if the blueshift is modeled only by multiple scattering theories, or if effective medium theories can reproduce such spectral behavior.
\end{itemize}


